\chapter{Black Holes as Zero-Entropy States}

\section{The Method}

In the previous chapter, we established a radical equivalence:
a physical system and the static data trace of its simulation are not distinct entities locked in a causal chain.
They are two representationally equivalent sides of the exact same mathematical coin.
When we look at a physical process, we are looking at geometry; when we look at its execution trace, we are looking at code. 

This realization provides us with a physical tool.
When our geometric descriptions hit a brick wall, we do not have to stop.
We can step through the looking glass, look at the code representation,
and read the answers back into the physics.


\section{The Problem}

There is no place in the cosmos where geometry fails more spectacularly
than at the center of a black hole: the classical singularity. 

Standard General Relativity dictates that when a massive star collapses under its own gravity,
its matter packs into an infinitely dense, infinitely small point where spacetime curvature becomes infinite.
This is where physics supposedly breaks down. 

To track what happens during gravitational collapse without relying on diverging geometric coordinates,
we translate the physical state of a collapsing system into a raw data stream. 

Imagine a large cloud of physical particles collapsing inward.
We discretize their positions, map their coordinates to fixed-width binary integers,
and parse the resulting system state as a raw bitstring.
We then observe this bitstring through three distinct informational lenses:

\begin{enumerate}
\item \textbf{Single-Bit Shannon Entropy:} This measures the fundamental diversity of the data.
  It counts the balance of zeros and ones across the machine's memory.
  If a system is highly ordered or completely uniform, this value drops.
\item \textbf{Pattern Block Entropy:} This groups the bits into consecutive short words (e.g., $4$-bit sequences)
  to check for localized correlations.
  It tells us if the data is forming repeating computational patterns.
  \textbf{Dynamic Spectral Complexity:} This isolates the structural variations of the cloud from its bulk movement.
  By computing a discrete Fourier transform on the spatial deviations of the particles, it tracks how many distinct
  spatial frequencies are required to render the shape of the cloud.
  If the cloud stretches, distorts, or scrambles, this metric spikes. If the cloud becomes perfectly uniform,
  it falls to zero.
\end{enumerate}

When simulating this collapse across different black hole geometries—such as a static
Schwarzschild black hole or a rapidly spinning Kerr black hole—these informational metrics reveal a
clear picture of what a singularity actually is.

\section{Inside the Event Horizon}

As the simulated particle cloud crosses the event horizon and heads toward the center,
the bitwise Shannon entropy and pattern block entropy decline smoothly and monotonically across
all geometries. The configuration space is shrinking; data variety is actively bleeding out of the system.

The simulation will crash due to numeric instabilities before the singularity forms. However, this
does not matter as we can extrapolate the entropy curves and draw the conclusion; they point towards
the same point; zero entropy state.


\section{The Mathematical Proof of a Point}

What does it mean when the entropy of a data trace becomes exactly zero? 

In computer science, a bitstring has zero Shannon entropy if and only if every
single bit in the system is identical—either a solid string of zeros ($0000\dots$) or a
solid string of ones ($1111\dots$). 

Now, let us pass that zero-entropy bitstring back through our decoding map to turn it back into physical geometry.
What does a solid string of identical bits decode to? It decodes to a single, identical coordinate triple. 

No matter what coordinate system you choose, no matter what simulation architecture you use,
and no matter how you define your decoding map, an information stream with zero variety can
only ever decode to a single, structureless geometric object: a point of zero size.

This yields three profound insights for physics:

\begin{itemize}
\item \textbf{The Singularity is Not Infinite:} General Relativity breaks down because its
  geometric equations try to calculate curvature over a manifold.
  But a manifold requires a non-trivial bitstring to map distinct, neighboring points.
  A zero-entropy state provides no variety. Spacetime breaks down not because the physics has
  exploded to infinity, but because there is nothing left to describe. 
\item \textbf{The Absence of Matter:} A zero-entropy point possesses no internal structural capacity.
  It cannot host microstructures, it cannot differentiate coordinates, and it cannot contain individual particles. 
\item \textbf{The Core is Empty:} Because a zero-entropy state cannot hold particles
  or microstates, it is difficult to see how it could support stress-energy tensor at its coordinate.
  With no matter present at the point itself, there is no infinite mass source at the center
  causing an infinite gravitational field. The singularity is completely empty.
\end{itemize}

\section{No New Physics Required}

Modern theoretical physics expends immense effort trying to ``fix'' the singularity.
Loop Quantum Gravity attempts to quantize spacetime at the Planck scale to create a cosmic ``bounce.''
String Theory tries to replace the center with a dense, horizonless ``fuzzball.''
Both approaches are forced to invent entirely new, unverified physics to smooth out the mathematical wrinkles of infinity.

This framework of Substrate Independence is appealing because it requires no new physics.
We do not need to invent a quantum gravity regularizer to halt the collapse before it hits infinity.
The regularization is baked directly into the informational structure of the universe. 

When a physical system collapses to absolute zero entropy, the geometry automatically becomes a trivial,
structureless point. The singularity is resolved not by stopping it, but by recognizing that it is the simplest,
most harmless shape a data structure can take.

\section{Holographic Reconciliation}

This informational descent to zero entropy solves an apparent paradox regarding black hole thermodynamics. 

According to the Bekenstein-Hawking formula, the entropy of a black hole is proportional to the surface area
of its event horizon. As a star collapses, the horizon grows, meaning the black hole's entropy increases.
How can the exterior entropy grow while our execution trace shows the interior geometry collapsing to zero entropy?

They are measuring two completely different things. The Bekenstein-Hawking entropy measures the total
informational capacity accessible to an outside observer looking at the boundary.
The execution-trace entropy measures the localized, internal geometric variety of the matter inside. 

As matter falls past the horizon, the information describing its internal configurations is stripped away
from the interior coordinate space, driving its local entropy to zero, while its informational signature is
conserved holographically on the boundary. The two descriptions are perfectly consistent.

\section{The Takeaway}

By utilizing the execution trace as a rigorous physical proxy,
we have demystified the most terrifying obstacle in classical gravity.
The singularity is not a physical catastrophe where the laws of nature shatter into infinite density.
It is merely an informational boundary condition—the natural limit where the variety of a
spatial system is completely compressed. 


